\subsection{Proof of Correctness}

\noindent\textbf{Theorem: Optimality of Greedy Green-Light Window Scheduling}

Let $W = \{w_1, w_2, \ldots, w_n\}$ represent a set
of candidate green-light windows. Each window
$w_i$ has a starting time $s_i$ and finishing time
$f_i$, where $s_i < f_i$.

The earliest-finish-time greedy algorithm produces
a maximum-cardinality subset of mutually compatible
green-light windows.

\medskip
\noindent\textbf{Proof.}

Let $g$ be the window with the earliest finishing
time, and let $O = (o_1, o_2, \ldots, o_k)$ be an
optimal schedule containing the maximum possible
number of compatible windows.

Since the greedy algorithm selects the window
with the earliest finishing time,

\[
f(g) \leq f(o_1).
\]

We can replace $o_1$ with $g$ in the optimal
schedule without reducing the number of
selected windows.

This replacement is valid because every
subsequent window in $O$ starts at or after
$f(o_1)$, which is greater than or equal to $f(g)$.

Therefore, the modified schedule remains
compatible and contains the same number
of windows as the original optimal schedule.

After selecting $g$, the algorithm considers
only windows beginning at or after $f(g)$.

The remaining problem has the same structure
as the original scheduling problem.

By repeatedly applying the same exchange argument
to the remaining windows, every greedy choice
can be included in an optimal schedule.

Therefore, the earliest-finish-time greedy
algorithm always produces a maximum-cardinality
set of mutually compatible green-light windows.

\hfill $\square$